\documentclass[11pt]{article}
\usepackage[a4paper,margin=26mm]{geometry}
\usepackage{amsmath,amssymb,amsthm,booktabs,hyperref}
\newcommand{\E}{\mathbb E}\newcommand{\Pp}{\mathbb P}
\newtheorem{theorem}{Theorem}\newtheorem{lemma}[theorem]{Lemma}
\title{Inward displacement pays for conditional noise}
\author{}\date{28 October 2026}
\begin{document}\maketitle

A centered noise cannot have large variance and a small upper endpoint
without a negative third moment. Keeping that third moment in the gap
budget gives a stronger noise cost. This argument uses the actual central
gap; it needs neither the outer support strips nor a count comparison.
It reduces the acquired global noise ceiling from $.616$ to $.36$, and
to $.32$ on the left range. It also confines the part below $p=1/3$ to
the indices $42\le n\le168$.

\section{The common-law geometric inequality}
Let $X$ be centered, with variance one and compact support, with an
actual central gap $(-A_-,A_+)$ and no atom at zero.
Let $n$ be a positive integer. Write
\[
 p=\Pp(X>0),\quad q=1-p,\quad v=pq,\quad
 m=\E X_+=\E X_-,\quad\ell=m/v.
\]
Use its actual sign $B=\mathbf1_{X>0}$ to write
$X=\ell(B+U-p)$, so $\E(U\mid B)=0$. Put
\[
 a_0=p-A_-/\ell,\qquad a_1=q-A_+/\ell,\qquad
 d=a_0+a_1,\quad s_j=\E(U^2\mid B=j),
\]
\[
 s=qs_0+ps_1,\quad\lambda=ns,\quad
 \lambda_0=nqs_0,\quad\lambda_1=nps_1.
\]
Thus $a_j\ge0$, $U_0\le a_0$, $U_1\ge-a_1$,
$\lambda_0+\lambda_1=\lambda$, and $\ell=(v+s)^{-1/2}$.
For $\sigma=\sqrt{10}$ define
\[
 Q(p)=(\sigma+1)p^2+(\sigma-1)q^2-6pq
       =(6+2\sigma)(p-p_*)^2,\qquad p_*=(4-\sigma)/2,
\]
\[
 D=(3+\sigma)\rho-\gamma-3J,\quad
 \rho=\E|X|^3,\quad\gamma=\E X^3,
 \quad J=\sum_{(a,b)}\tau_{a,b}(b-a)^2.
\]
The sum is over the bounded support gaps and
$\tau(t)=\E[X\mathbf1_{X\ge t}]$ is constant on each gap.
There are at most countably many gaps. Their nonnegative contributions
are interpreted as a monotone sum; the integral bounds below apply to
each finite subcollection and then to the full sum by monotone convergence.
Compact support makes the displayed moment bounds finite.
All quantities in this note refer to one actual law.

\begin{lemma}[One-sided cubic cost]\label{g:cubic}
If $Z$ is centered, $Z\le a$ almost surely, $a>0$, and
$t=\E Z^2$, then
\[
                 \E Z^3\le at-t^2/a.
\]
If $a=0$, centering forces $Z=0$ almost surely.
\end{lemma}
\begin{proof}
Expand the nonnegative expectation
$\E[(a-Z)(Z+t/a)^2]=at-\E Z^3-t^2/a$.
No lower support endpoint is used. Apply this also to $-U_1$.
\end{proof}

Fix $0\le\bar d<1$ with $d\le\bar d$, and define
\begin{align*}
 Y_0(p)&=3(\sigma+3)p-(\sigma+7)\bar d
       +2\sqrt{6(1-\bar d)(\sigma+1)p},\\
 Y_1(p)&=3(\sigma+1)q-(\sigma+5)\bar d
       +2\sqrt{6(1-\bar d)(\sigma-1)q}.
\end{align*}

\begin{theorem}[A probability--noise--allocation budget]\label{g:budget}
Under the preceding hypotheses,
\begin{equation}\label{g:newbudget}
 nD\ge\frac{nvQ(p)+\lambda_0Y_0(p)+\lambda_1Y_1(p)}
                    {(v+\lambda/n)^{3/2}}.
\end{equation}
In particular the numerator is at least
$nvQ(p)+\lambda\min(Y_0(p),Y_1(p))$.
\end{theorem}
\begin{proof}
We retain the two residual noise measures separately in the direct gap
budget. For example the positive outer gaps satisfy
\[
 J_o^+\le2\int_{A_+}^{\infty}(t-A_+)\tau(t)\,dt
       =\E[X_+(X_+-A_+)^2],
\]
with the reflected expression on the negative side. Adding the actual
central gap gives the exact lower-budget expansion used in v10:
\begin{align*}
 nD(v+s)^{3/2}\ge{}&nvQ+6nv\{d(1-d)+a_0a_1\}\\
 &+3(\sigma+3)p\lambda_0+3(\sigma+1)q\lambda_1
       -6a_0\lambda_0-6a_1\lambda_1\\
 &-(\sigma+1)nq\E U_0^3+(\sigma-1)np\E U_1^3.
\end{align*}
This is the expansion of the common geometric cost before replacing a
conditional cubic by its endpoint bound; see
\path{mathematics/simplification-v10/2026-10-28/geometry.tex},
the proof of the common geometric cost. The integration includes endpoint atoms;
individual integration endpoints have Lebesgue measure zero.

Lemma~\ref{g:cubic} improves the last line, giving
\begin{align*}
 nD(v+s)^{3/2}\ge{}&nvQ+6nv\{d(1-d)+a_0a_1\}\\
 &+\{3(\sigma+3)p-(\sigma+7)a_0\}\lambda_0
  +\{3(\sigma+1)q-(\sigma+5)a_1\}\lambda_1\\
 &+\frac{(\sigma+1)\lambda_0^2}{nqa_0}
  +\frac{(\sigma-1)\lambda_1^2}{npa_1}.
\end{align*}
A term with $a_j=0$ has $\lambda_j=0$ and is simply absent;
there is no division by zero. Bound $a_j\le\bar d$ in the two negative
linear terms and use
$d(1-d)+a_0a_1\ge(1-\bar d)(a_0+a_1)$.
For each $j$, apply the arithmetic--geometric mean inequality to its
central-displacement and reciprocal-cubic costs. For example
\[
 6nv(1-\bar d)a_0+
       \frac{(\sigma+1)\lambda_0^2}{nqa_0}
 \ge2\lambda_0\sqrt{6(1-\bar d)(\sigma+1)p}.
\]
The other group is identical with $p,q$ and $\sigma+1,\sigma-1$
interchanged. This proves \eqref{g:newbudget}.
\end{proof}

The allocation $\lambda_1/\lambda$ has not been optimized independently
of the law. The full inequality may be retained as a shared resource;
replacing its affine allocation term by its smaller endpoint is only a
conservative simplification. No contact, convolution or predecessor
assumption was used in Theorem~\ref{g:budget}.

\section{Consequences for the acquired domain}
The following are existing global geometric outputs of the selected
coarse entry, not conclusions of a count calculation:
\begin{equation}\label{g:input}
 n\ge42,\quad p\in[1/4,27/50],\quad d\le17/250,
 \quad0\le\lambda\le77/125,\quad nD\le B_*:=32.20665378.
\end{equation}
The strict $p>1/4$, $d<17/250$ and $\lambda<77/125$ previously proved
may safely be weakened to these closed premises. Only the incoming
coarse entry and direct gap geometry are needed; the $.6/ .7$ support
strips are not needed for the conclusions below.

\begin{theorem}\label{g:ceiling}
Under \eqref{g:input}, one has the following strict noise ceilings:
\[
\begin{array}{c|cccc}
 p\text{ at most}&3/10&1/3&2/5&27/50\\\hline
 \lambda\text{ less than}&1/6&1/4&8/25&9/25
\end{array}
\]
Moreover the following integer-index implications hold:
\[
\begin{array}{c|rrrrrr}
 p\text{ at most}&.30&1/3&.35&.375&.39&.40\\\hline
 n\text{ at most}&84&168&262&657&1530&3598.
\end{array}
\]
\end{theorem}
\begin{proof}
Set $H_j=Y_j/v$. On the stated domain,
$v\ge3/16$, $0\le Q\le9/25$ and $H_j>30$.
These elementary bounds are checked with rational square comparisons.
For fixed $p$, put
\[
 F_j(n,\lambda)=
 \frac{nQ+H_j\lambda}{\sqrt v\{1+\lambda/(nv)\}^{3/2}}.
\]
Theorem~\ref{g:budget} gives $nD\ge\min_jF_j$.
Writing $z=\lambda/n\le11/750$, the sign of its derivative in
$\lambda$ is that of
\[
 H_j\{1-z/(2v)\}-3Q/(2v)
 \ge30\{1-(11/750)/(2(3/16))\}
               -3(9/25)/(2(3/16))>0.
\]
For fixed positive $\lambda$, $F_j$ also increases in $n$:
apart from a positive factor it is
$(Q/z+H_j)/(1+z/v)^{3/2}$, which decreases in $z$.
Thus to exclude $\lambda\ge L$ it suffices to check
$F_0(42,L)>B_*$ and $F_1(42,L)>B_*$ on the relevant probability
interval. This is an algebraic monotonicity argument, not a perturbation
or limit of actual laws.

Here is a small exact certificate for these eight continuous inequalities.
Use
\[s_-=3.162277660<\sigma<s_+=3.162277661.\]
In each $Y_j$ replace its positive coefficients by $s_-$ and its negative
coefficient by $s_+$. Replace its square-root term by the chord joining
rational lower bounds at the two interval endpoints. Concavity of the
square root puts this affine chord below the square root on the entire
interval. Call the resulting affine polynomial $y_j(p)$, and let $Q_-$
denote $Q$ with $\sigma$ replaced by $s_-$.
For $N=42$ form
\[
 R_j(p)=NvQ_-(p)+Ly_j(p),\qquad
 P_j(p)=N^3R_j(p)^2-B_*^2(Nv+L)^3.
\]
Positive $R_j$ and $P_j$ prove the desired inequality without a sign
ambiguity from squaring. The polynomial degrees are four and eight.
For the four intervals and ceilings in the first table, all degree-four
Bernstein coefficients are positive on the whole interval. The degree-eight
coefficients are also positive on the whole interval, except for the global
$j=1$ polynomial, for which the four equal closed subintervals suffice.
There are 139 strictly positive coefficients in total, on eleven closed
polynomial intervals. The script constructs the polynomials from these
formulas, rather than loading fitted coefficients. It also verifies each
square-root enclosure before using it. This proves the first table.

Monotonicity at $\lambda=0$ gives $D\ge Q(p)/\sqrt{pq}$.
On $[1/4,2/5]$ this ratio decreases: $Q$ is positive and decreasing,
while $pq$ is increasing. For each upper endpoint $b$ and index $N_b$
in the second table, the exact positive-base comparison is
\[
             (N_b+1)^2Q_-(b)^2>B_*^2b(1-b).
\]
It excludes $n\ge N_b+1$ when $p\le b$, proving the second table.
\end{proof}

The same inequality gives explicit near-critical localization:
\[
 |p-p_*|\le
 \sqrt{\frac{B_*}{2(6+2\sqrt{10})n}}
 <\frac{23}{20\sqrt n}.
\]
In particular no acquired witness with $n\ge169$ lies below $1/3$,
and no acquired witness with $n\ge3599$ lies at or below $.4$.
These remove whole probability/index regions before count comparison;
they do not certify the remaining finite indices.

\section{The Bernoulli gap: exactly the available scope}
Schulz's original thesis, \emph{The optimal Berry--Esseen constant in the
binomial case} (2016), Theorem~1 on printed page~1, supplies the strict
finite-$n$ bound at $C_E$ for every $0<p<1$. Its sharper formula is stated
only for $p\in[1/3,2/3]$; see the
\href{https://d-nb.info/1197702695/34}{original thesis}.
For $1/3\le p\le1/2$, that formula yields the weighted reserve
\[
 C_E\rho_B-\sqrt n\,\sup_x|F_B(x)-\Phi((x-np)/\sqrt{nv})|
 >\frac{Q(p)}{6\sqrt{2\pi v}}
   =\frac{2C_E(p-p_*)^2}{\sqrt v}.
\]
This is an exact algebraic identity, not a new extension of the external
theorem. Below $1/3$ we retain only strictness. The geometric index bound
now confines that unresolved quantitative-gap range to $42\le n\le168$.
No Bernoulli-gap premise is used in the new noise ceiling.

The thesis gives a warning example on printed pages 26--27:
$n=217$, $p=77/2604$, evaluated at $k=6$. Our exact rational normal
enclosure independently confirms that the sharper formula fails there.
This lies outside the acquired probability interval and does not refute
its possible extension to $[1/4,1/3)$; it rules out an unrestricted
extrapolation.

\section{Checks, an obstruction, and the remaining work}
The scalar checker uses exact rational arithmetic for the 139 coefficients
and all denominator, monotonicity and integer-index gates. A separate
actual-law check computes every support gap and its weight, retaining
threshold atoms, rather than treating conditional means as support points.
It includes zero conditional noise, both variance allocations, symmetric
and skew Bernoulli limits, tiny positive defects, and inward cubic equality
cases. Each positive-noise law is also checked after reflection, with its
actual sign probability and conditional variance allocations exchanged.
These examples test the formulas; they do not prove the continuous
domain by sampling.

There is a concrete limit to the reduction. Take $n=42$, $p=3/10$,
$U_1=0$, and let the centered law $U_0$ be supported on
$\{-3/5,5/882\}$. Then $\lambda=1/10$, $d=5/882$,
and the standardized law obeys the global support and geometric bounds,
including $nD<27<32.20665378$. It lies below $1/3$ and above the completed
left count range $.01$. Thus these necessary geometric data alone do not
remove every remaining left case. No contact or nonnegative-excess premise
is asserted for this example.

The new bounds make all count work at $\lambda\ge.36$ unnecessary for
an acquired witness, and all left count work at $\lambda\ge.32$
unnecessary. Smaller left probability ranges have the stronger ceilings
in Theorem~\ref{g:ceiling}. Their derivation requires no outer-strip or
heat/ray calculation. The coarse entry, actual gap, original geometric
bootstrap and their contact/predecessor ancestry remain. Count comparison
below the new ceilings, signed/reflected acquisition, and the independent
finite base through 41 remain open obligations. No previous theorem,
certificate or manuscript source is changed.
\end{document}
